% Chapter 7: Implementing Off-Chain Logic for ZK in Cardano.
% All source code lives in zk-password/ at the root of the repository.

%--------------------------------------------------------------------------
\section{The On-Chain Part Is the Small Part} \label{sec:ch7-small-part}

Chapter~\ref{ch:password} produced a working ZK validator. Let us look at how
big it actually is. After \texttt{aiken build}, the compiled validator,
including the entire Groth16 verifier, the pairing equation, and the public
input accumulation, occupies \textbf{570 bytes} of Plutus Core.

Everything else in a ZK application is off-chain: the circuit, the trusted
setup, the witness calculation, the prover, the point compression, the datum
and redeemer encoding, the transaction construction, the UTxO indexing, and the
key management. That is thousands of lines of code and several hundred megabytes
of artifacts, against 570 bytes on-chain.

\begin{highlight}
A ZK application on Cardano is an off-chain application with a 570-byte
on-chain referee. If you get the off-chain pipeline wrong, nothing works,
and the on-chain error message will only ever tell you \emph{``validation
failed''}.
\end{highlight}

This chapter is about that off-chain pipeline: which software computes proofs,
how to drive it step by step, how to get a proof into a redeemer with MeshJS,
and what the resulting transactions actually look like and cost.

\subsection*{Division of Responsibilities}

\begin{center}
\begin{tabular}{lll}
  \hline
  \textbf{Responsibility}       & \textbf{Where}  & \textbf{Tool} \\
  \hline
  Define the statement          & Off-chain       & Circom \\
  Produce public parameters     & Off-chain       & snarkjs (ceremony) \\
  Compute the witness           & Off-chain       & WASM / C++ \\
  Produce the proof             & Off-chain       & snarkjs / rapidsnark \\
  Compress curve points         & Off-chain       & @noble/curves \\
  Encode as Plutus Data         & Off-chain       & MeshJS \\
  Build \& submit transaction   & Off-chain       & MeshJS + provider \\
  \textbf{Verify the proof}     & \textbf{On-chain} & \textbf{Aiken} \\
  Enforce spending conditions   & On-chain        & Aiken \\
  \hline
\end{tabular}
\end{center}

Note that only two rows are on-chain, and the cryptographic one is a single
function call. Cardano's contribution to a ZK application is remarkably
narrow, and, as we will measure below, remarkably cheap.

%--------------------------------------------------------------------------
\section{What Software Computes Proofs} \label{sec:ch7-provers}

There is no single ``ZK prover''. There is a family of tools, and the right
choice depends on where the proof must be computed: in a browser tab, on a
server, or on a phone.

\subsection*{The Circuit Compiler}

\textbf{Circom} is the constraint-system compiler we have been using. It takes
\texttt{.circom} source and emits an R1CS file plus a witness calculator. Circom
does not prove anything; it only describes the statement and computes witnesses.

\begin{verbatim}
circom --r1cs --wasm --sym --inspect \
       --prime bls12381 circuit/circuit.circom -o circuit/
\end{verbatim}

The \texttt{--prime bls12381} flag is not optional for Cardano. Circom's default
is BN254 (Ethereum's curve). A BN254 proof is mathematically valid and will
verify with snarkjs, and will fail on Cardano, because Cardano's builtins
only implement BLS12-381. This mismatch is the single most common cause of
``my proof verifies locally but the validator rejects it''.

Alternatives to Circom exist and are worth knowing about, though none of them
currently has a Cardano verifier:

\begin{itemize}
  \item \textbf{Noir} (Aztec): a Rust-like ZK language, much nicer to write
        than Circom, but its default backend (Barretenberg / UltraHonk) has no
        Aiken verifier.
  \item \textbf{Halo2} (Zcash / PSE): no trusted setup, recursive-friendly,
        but its verifier needs operations Cardano does not expose cheaply.
  \item \textbf{Circom + PLONK}: snarkjs can produce PLONK proofs from the
        same R1CS. The verifier is more complex than Groth16 and, at the time
        of writing, no production-grade PLONK verifier exists in Aiken.
        We return to this in Chapter~\ref{ch:future}.
\end{itemize}

\begin{remark}
For Cardano today, \textbf{Circom + snarkjs + Groth16 over BLS12-381} is the
only fully supported path from circuit to on-chain verification. Everything in
this chapter assumes it.
\end{remark}

\subsection*{The Provers}

Given an R1CS, a proving key, and a witness, several programs can produce the
proof. They differ by orders of magnitude in speed.

\begin{center}
\begin{tabular}{llll}
  \hline
  \textbf{Prover} & \textbf{Language} & \textbf{Runs in} & \textbf{Speed} \\
  \hline
  snarkjs            & JavaScript / WASM & Browser, Node & Baseline \\
  rapidsnark         & C++ (asm)         & Server, native & 5--10$\times$ \\
  ark-circom         & Rust (arkworks)   & Server, native & 5--10$\times$ \\
  witnesscalc        & C++               & iOS, Android   & 10$\times$ (witness) \\
  gnark              & Go                & Server         & 5--10$\times$ \\
  \hline
\end{tabular}
\end{center}

\begin{itemize}
  \item \textbf{snarkjs} is the reference implementation and the one to start
        with. It runs everywhere, needs no compilation, and is what we use in
        this book. Install with \texttt{npm install -g snarkjs}.
  \item \textbf{rapidsnark} is the production choice for server-side proving.
        It reads the same \texttt{.zkey} and \texttt{.wtns} files as snarkjs and
        emits an identical \texttt{proof.json}, so it is a drop-in replacement
        once your pipeline works. Note that rapidsnark's BLS12-381 support is
        less mature than its BN254 support, verify this before committing.
  \item \textbf{ark-circom} lets you drive Circom circuits from Rust using
        arkworks. Useful if your off-chain service is already Rust, and
        arkworks has first-class BLS12-381 support.
  \item \textbf{witnesscalc} replaces only the witness-calculation step, which
        is often the bottleneck on mobile hardware.
\end{itemize}

\subsection*{The Supporting Cast}

\begin{center}
\begin{tabular}{ll}
  \hline
  \textbf{Package} & \textbf{Role} \\
  \hline
  \texttt{circomlib}       & Circuit primitives: Poseidon, comparators, bits \\
  \texttt{circomlibjs}     & JS reimplementation of those primitives \\
  \texttt{@noble/curves}   & Point compression to Cardano's byte format \\
  \texttt{@meshsdk/core}   & Transaction building, Plutus Data encoding \\
  \texttt{@lucid-evolution/lucid} & Alternative to MeshJS \\
  Blockfrost / Koios / Maestro & Chain queries and submission \\
  \hline
\end{tabular}
\end{center}

\begin{remark}
\texttt{circomlibjs} is a trap for Cardano developers. Its \texttt{buildPoseidon}
implements Poseidon over \textbf{BN254}, not BLS12-381. Calling it to compute the
hash that your BLS12-381 circuit will check produces a value that is simply
wrong, and the failure surfaces much later as an unsatisfiable constraint.
Chapter~\ref{ch:password} solved this by compiling a tiny helper circuit
(\texttt{poseidon\_hash.circom}) and reading the hash out of its witness.
Do that.
\end{remark}

%--------------------------------------------------------------------------
\section{The Pipeline, Step by Step} \label{sec:ch7-pipeline}

The full pipeline has six stages. Stages 1--3 happen once per circuit. Stages
4--6 happen once per proof. Confusing the two is the second most common source
of bugs.

\begin{center}
\begin{tabular}{clll}
  \hline
  \textbf{\#} & \textbf{Stage} & \textbf{Frequency} & \textbf{Output} \\
  \hline
  1 & Compile circuit    & Per circuit & \texttt{.r1cs}, \texttt{.wasm} \\
  2 & Trusted setup      & Per circuit & \texttt{.zkey}, \texttt{vk.json} \\
  3 & Compress the VK    & Per circuit & \texttt{compressed\_vk.json} \\
  4 & Compute witness    & Per proof   & \texttt{.wtns} \\
  5 & Prove              & Per proof   & \texttt{proof.json}, \texttt{public.json} \\
  6 & Compress the proof & Per proof   & \texttt{compressed\_proof.json} \\
  \hline
\end{tabular}
\end{center}

\subsection*{Stage 1: Compile}

\begin{verbatim}
circom --r1cs --wasm --sym --inspect \
       --prime bls12381 circuit/circuit.circom -o circuit/

snarkjs r1cs info circuit/circuit.r1cs
\end{verbatim}

Read the \texttt{r1cs info} output carefully every time. It tells you the curve,
the number of public inputs, and the number of constraints. The number of public
inputs must equal \texttt{length(vk\_ic) - 1} in your compressed verification
key, and must equal the length of the list you pass to
\texttt{groth16\_verify} in Aiken. For our password circuit:

\begin{verbatim}
[INFO]  snarkJS: Curve: bls12-381
[INFO]  snarkJS: # of Wires: 219
[INFO]  snarkJS: # of Constraints: 216
[INFO]  snarkJS: # of Public Inputs: 1
[INFO]  snarkJS: # of Private Inputs: 1
[INFO]  snarkJS: # of Outputs: 0
\end{verbatim}

\subsection*{Stage 2: Trusted Setup}

Phase 1 (Powers of Tau) is circuit-independent and reusable. Phase 2 is bound to
this specific R1CS.

\begin{verbatim}
# Phase 1: universal, reusable across circuits
snarkjs powersoftau new bls12381 12 pot12_0000.ptau -v
snarkjs powersoftau contribute pot12_0000.ptau pot12_0001.ptau \
    --name="dev" -e="some entropy"
snarkjs powersoftau beacon pot12_0001.ptau pot12_beacon.ptau \
    00000000000000000000000000000000 10 -n="Dev beacon"
snarkjs powersoftau prepare phase2 \
    pot12_beacon.ptau pot12_final.ptau -v

# Phase 2: specific to circuit.r1cs
snarkjs groth16 setup \
    circuit/circuit.r1cs pot12_final.ptau zkey_0000.zkey
snarkjs zkey contribute zkey_0000.zkey zkey_0001.zkey \
    --name="dev" -e="more entropy"
snarkjs zkey beacon zkey_0001.zkey zkey_final.zkey \
    00000000000000000000000000000000 10 -n="Final beacon"
snarkjs zkey export verificationkey \
    zkey_final.zkey verification_key.json
\end{verbatim}

The power \texttt{12} means $2^{12} = 4096$ maximum constraints. Choose the
smallest power that exceeds your constraint count: the \texttt{.ptau} and
\texttt{.zkey} file sizes grow linearly with $2^{\text{power}}$, and so does the
proving time. This matters enormously in practice; we will see in
Chapter~\ref{ch:tornado} that a circuit needing $2^{20}$ turns a 2\,MB
proving key into a multi-gigabyte one.

\begin{highlight}
\textbf{The proving key is not a secret, but it is not small either.}
Stage 2 is the point at which a ZK application acquires a distribution problem:
every prover needs \texttt{zkey\_final.zkey} and \texttt{circuit.wasm}, and
every verifier needs to agree on the same \texttt{compressed\_vk.json}.
\end{highlight}

\subsection*{Stage 3: Compress the Verification Key}

Cardano's BLS12-381 builtins accept only compressed points: 48 bytes for G1,
96 bytes for G2. snarkjs emits uncompressed affine coordinates. The
\texttt{compress\_vk.js} script from Chapter~\ref{ch:password} does the
conversion:

\begin{verbatim}
node circuit/compress_vk.js
\end{verbatim}

This is a per-circuit operation. The output is the value you will embed in the
datum (or hard-code in the validator, or store in a reference input, see
Section~\ref{sec:ch7-ops}).

\subsection*{Stage 4: Compute the Witness}

For our circuit the witness needs the password as a field element and its
Poseidon hash over BLS12-381:

\begin{verbatim}
node circuit/gen_input.js 123456
\end{verbatim}

\begin{verbatim}
password : 123456
hash     : 941543474650874718122864576426741288983580381
           5708024285868314342086739058356
written  : circuit/input.json
\end{verbatim}

\begin{remark}
\texttt{gen\_input.js} is the only place in the entire pipeline that touches the
plaintext password. Everything downstream operates on field elements and curve
points. When you audit a ZK application, this is the file to audit first.
\end{remark}

\subsection*{Stage 5: Prove}

\begin{verbatim}
snarkjs groth16 fullprove \
    circuit/input.json \
    circuit/circuit_js/circuit.wasm \
    zkey_final.zkey \
    proof.json public.json

snarkjs groth16 verify verification_key.json public.json proof.json
\end{verbatim}

Always run the off-chain \texttt{verify} before building a transaction. It costs
milliseconds and distinguishes ``my proof is wrong'' from ``my encoding is
wrong'', two failures that look identical from on-chain.

\subsection*{Stage 6: Compress the Proof}

\begin{verbatim}
node circuit/compress_proof.js
\end{verbatim}

The result is three hex strings. For our worked example:

\begin{verbatim}
{
  "pi_a": "b15f6da17265cbb596036075e87ddfe62d210a92c65cb30a
           c65375659ef26511c50b55b989abc6da22f05495fca74380",
  "pi_b": "ac6f06b66982fc793de1a1d393b66f80f967cf4f984bc18a
           c9e034a618d732f5a9dd9c632da14178b54584a682f80708
           08437cb5f80c0559f9cddf71f7edbddb4853249ba637ec3a
           59950eff969e99a2d4574a6f229a753cdb2d5674fcf060cf",
  "pi_c": "8d6da5b26036b096441e48b6c490f56c4a4f9148a881b0e5
           0c556626a4cef37d672c0e8d7c74fc0b2214e0bf63f36f8a"
}
\end{verbatim}

(Line breaks added for the page; the real file has one string per field.)
That is \textbf{192 bytes} of proof (48 + 96 + 48), regardless of how
complicated the statement was. This constant size is Groth16's defining
property and the reason it suits a blockchain.

%--------------------------------------------------------------------------
\section{Where Should the Prover Run?} \label{sec:ch7-where}

This is an architectural decision, not a technical one, and it determines
whether your application actually provides privacy. There are three options.

\subsection*{Architecture A: Client-Side Proving}

The witness and the proof are computed in the user's browser or app. The secret
never leaves the device.

\begin{verbatim}
  [user's browser]                          [chain]
  secret --> witness --> proof --------------> validator
             (wasm)      (snarkjs)             (Aiken)
\end{verbatim}

\begin{itemize}
  \item \textbf{Privacy:} maximal. Nobody but the user ever sees the secret.
  \item \textbf{Cost:} the user downloads the \texttt{.zkey} (megabytes to
        gigabytes) and pays the proving time.
  \item \textbf{Verdict:} the only architecture that honestly delivers
        zero-knowledge. Use it whenever the circuit is small enough.
\end{itemize}

\subsection*{Architecture B: Server-Side Proving}

The user sends the secret to a server, which computes the proof and returns it
(or builds the whole transaction). This is what Janus
Wallet\footnote{\texttt{https://github.com/leobel/janus-wallet}} does: a
Node.js service with \texttt{circom} installed, PostgreSQL for user and circuit
records, and Blockfrost for submission.

\begin{verbatim}
  [browser]        [prover service]           [chain]
  secret ----------> witness --> proof --------> validator
           (TLS)     (rapidsnark)
\end{verbatim}

\begin{itemize}
  \item \textbf{Privacy:} \emph{none against the server.} The server sees the
        password. The proof is still zero-knowledge with respect to the
        blockchain and every other observer, but the operator is fully trusted.
  \item \textbf{Cost:} fast proving, no large client download.
  \item \textbf{Verdict:} legitimate for onboarding and UX, provided you say so
        plainly. It is a custodial-style trust assumption wearing a
        zero-knowledge badge.
\end{itemize}

\begin{highlight}
``The proof is zero-knowledge'' and ``the system is zero-knowledge'' are
different claims. A server-side prover satisfies the first and violates the
second. Be explicit about which one you are offering.
\end{highlight}

\subsection*{Architecture C: Hybrid}

The witness is computed client-side; the proof is computed server-side from the
witness. This does not help: for most circuits the witness contains the private
signals verbatim, so the server still learns the secret. The variant that
\emph{does} help is delegating only the transaction \emph{submission}, a
relayer. The relayer sees the proof (public anyway) and pays the fee, but never
sees the witness. We use exactly this pattern in Chapter~\ref{ch:tornado}.

%--------------------------------------------------------------------------
\section{Encoding Proofs as Plutus Data} \label{sec:ch7-encoding}

The bridge between snarkjs and Aiken is Plutus Data. Get the shape right and
everything works; get the field order wrong and you receive an unhelpful
validation failure.

\subsection*{The Mapping Rules}

\begin{center}
\begin{tabular}{lll}
  \hline
  \textbf{Aiken} & \textbf{Plutus Data} & \textbf{MeshJS} \\
  \hline
  \texttt{Int}            & \texttt{I} integer     & \texttt{number} / \texttt{bigint} \\
  \texttt{ByteArray}      & \texttt{B} bytes       & hex \texttt{string} \\
  \texttt{List<a>}        & \texttt{List}          & \texttt{Array} \\
  record / constructor    & \texttt{Constr} $n$    & \texttt{mConStr}$n$\texttt{([...])} \\
  \texttt{Option.Some(x)} & \texttt{Constr} 0      & \texttt{mConStr0([x])} \\
  \texttt{Option.None}    & \texttt{Constr} 1      & \texttt{mConStr1([])} \\
  \hline
\end{tabular}
\end{center}

Three rules govern the translation:

\begin{enumerate}
  \item \textbf{Field order is positional.} Aiken records serialise their fields
        in declaration order. The Plutus Data \texttt{Constr} carries no field
        names, so \texttt{mConStr0} arguments must appear in exactly the order
        the Aiken \texttt{type} declares them.
  \item \textbf{Single-constructor records use alternative 0.}
  \item \textbf{Opaque types cannot cross the boundary.} As we discovered in
        Chapter~\ref{ch:password}, \texttt{Scalar} is opaque; the datum must
        carry a plain \texttt{Int} and the validator converts it with
        \texttt{scalar.new}.
\end{enumerate}

\subsection*{Our Types, Encoded}

Recall the Aiken declarations:

\begin{verbatim}
pub type VerificationKey {
  alpha: ByteArray,        // G1: 48 bytes
  beta:  ByteArray,        // G2: 96 bytes
  gamma: ByteArray,        // G2: 96 bytes
  delta: ByteArray,        // G2: 96 bytes
  ic:    List<ByteArray>,  // G1 points
}

pub type Proof {
  pi_a: ByteArray,  pi_b: ByteArray,  pi_c: ByteArray,
}

pub type Datum    { credential_hash: Int, vk: VerificationKey }
pub type Redeemer { proof: Proof }
\end{verbatim}

And the corresponding MeshJS encoding:

\begin{verbatim}
import { mConStr0 } from "@meshsdk/core";
import compressedVk from "./compressed_vk.json";
import compressedProof from "./compressed_proof.json";
import publicSignals from "./public.json";

// VerificationKey: five fields, in declaration order
const vkData = mConStr0([
  compressedVk.vk_alpha1,
  compressedVk.vk_beta2,
  compressedVk.vk_gamma2,
  compressedVk.vk_delta2,
  compressedVk.vk_ic,        // List<ByteArray> -> JS Array
]);

// Datum { credential_hash: Int, vk: VerificationKey }
// public.json holds a decimal string; it exceeds 2^53, so BigInt
// is mandatory. parseInt() here silently corrupts the value.
const datumData = mConStr0([
  BigInt(publicSignals[0]),
  vkData,
]);

// Proof, then Redeemer { proof: Proof }
const proofData = mConStr0([
  compressedProof.pi_a,
  compressedProof.pi_b,
  compressedProof.pi_c,
]);

const redeemerData = mConStr0([proofData]);
\end{verbatim}

\begin{remark}
The \texttt{credential\_hash} is a 255-bit field element. JavaScript's
\texttt{number} holds 53 bits of integer precision. Using \texttt{parseInt} or
\texttt{Number} on it produces a nearby-but-wrong value, the datum hash changes,
and the validator fails on a value that \emph{looks} right in the console.
Always \texttt{BigInt}.
\end{remark}

\subsection*{Sanity-Check the Encoding Before Submitting}

You can verify the datum round-trips without spending anything:

\begin{verbatim}
import { serializeData, deserializeDatum } from "@meshsdk/core";

const cbor = serializeData(datumData, "Mesh");
console.log("datum cbor:", cbor);
console.log("round-trip:", deserializeDatum(cbor));
\end{verbatim}

Compare the CBOR against what \texttt{aiken blueprint} reports for the datum
schema. A mismatch here is a mismatch on-chain.

%--------------------------------------------------------------------------
\section{Building the Transactions with MeshJS} \label{sec:ch7-mesh}

We now have all the pieces. Let us build the two transactions: one to lock ADA
behind the ZK lock, one to unlock it with a proof.

\subsection*{Setup: \texttt{common.ts}}

\begin{verbatim}
npm install @meshsdk/core @meshsdk/core-csl
\end{verbatim}

\begin{verbatim}
import fs from "node:fs";
import {
  BlockfrostProvider, MeshTxBuilder, MeshWallet,
  serializePlutusScript, UTxO,
} from "@meshsdk/core";
import { applyParamsToScript } from "@meshsdk/core-csl";
import blueprint from "../plutus.json";

const projectId = process.env.BLOCKFROST_PROJECT_ID;
if (!projectId) throw new Error("Missing BLOCKFROST_PROJECT_ID");

export const provider = new BlockfrostProvider(projectId);

export const wallet = new MeshWallet({
  networkId: 0,             // 0 = preprod/preview, 1 = mainnet
  fetcher: provider,
  submitter: provider,
  key: {
    type: "root",
    bech32: fs.readFileSync("me.sk").toString(),
  },
});

export function getScript() {
  const scriptCbor = applyParamsToScript(
    blueprint.validators[0].compiledCode, []
  );
  const scriptAddr = serializePlutusScript(
    { code: scriptCbor, version: "V3" }
  ).address;
  return { scriptCbor, scriptAddr };
}

export function getTxBuilder() {
  return new MeshTxBuilder({
    fetcher: provider,
    submitter: provider,
    evaluator: provider,
  });
}

export async function getUtxoByTxHash(
  txHash: string
): Promise<UTxO> {
  const utxos = await provider.fetchUTxOs(txHash);
  if (utxos.length === 0) throw new Error("UTxO not found");
  return utxos[0];
}
\end{verbatim}

Passing \texttt{evaluator} lets MeshJS ask the provider to evaluate the script
and fill in the execution-unit budget automatically, instead of you guessing it.

\subsection*{\texttt{lock.ts}: Locking Behind the ZK Lock}

\begin{verbatim}
import { Asset, mConStr0 } from "@meshsdk/core";
import { getScript, getTxBuilder, wallet } from "./common";
import compressedVk from "../compressed_vk.json";
import publicSignals from "../public.json";

async function main() {
  const assets: Asset[] =
    [{ unit: "lovelace", quantity: "5000000" }];

  const utxos = await wallet.getUtxos();
  const walletAddress = (await wallet.getUsedAddresses())[0];
  const { scriptAddr } = getScript();

  const vkData = mConStr0([
    compressedVk.vk_alpha1, compressedVk.vk_beta2,
    compressedVk.vk_gamma2, compressedVk.vk_delta2,
    compressedVk.vk_ic,
  ]);

  const datum = mConStr0([BigInt(publicSignals[0]), vkData]);

  const txBuilder = getTxBuilder();
  await txBuilder
    .txOut(scriptAddr, assets)
    .txOutInlineDatumValue(datum)
    .changeAddress(walletAddress)
    .selectUtxosFrom(utxos)
    .complete();

  const signedTx = await wallet.signTx(txBuilder.txHex);
  const txHash = await wallet.submitTx(signedTx);
  console.log("locked at:", txHash);
}

main();
\end{verbatim}

We use \texttt{txOutInlineDatumValue}, not \texttt{txOutDatumHashValue}. With an
inline datum the verification key and the Poseidon hash live directly in the
UTxO, so any spender can read them from the chain without a side channel. With a
datum hash the spender would have to obtain the preimage from somewhere else:
extra infrastructure for no privacy benefit, since neither value is secret.

\begin{highlight}
Storing the verification key in the datum makes each locked UTxO
self-describing, at the cost of roughly 430 bytes of on-chain data per UTxO.
Section~\ref{sec:ch7-ops} discusses the alternative: bake the VK into the
validator as a compile-time parameter.
\end{highlight}

\subsection*{\texttt{unlock.ts}: Spending with a Proof}

\begin{verbatim}
import { mConStr0 } from "@meshsdk/core";
import {
  getScript, getTxBuilder, getUtxoByTxHash, wallet,
} from "./common";
import compressedProof from "../compressed_proof.json";

async function main() {
  const lockTxHash = process.argv[2];

  const utxos = await wallet.getUtxos();
  const walletAddress = (await wallet.getUsedAddresses())[0];
  const collateral = (await wallet.getCollateral())[0];

  const { scriptCbor } = getScript();
  const scriptUtxo = await getUtxoByTxHash(lockTxHash);

  // Redeemer { proof: Proof { pi_a, pi_b, pi_c } }
  const redeemer = mConStr0([
    mConStr0([
      compressedProof.pi_a,
      compressedProof.pi_b,
      compressedProof.pi_c,
    ]),
  ]);

  const txBuilder = getTxBuilder();
  await txBuilder
    .spendingPlutusScript("V3")
    .txIn(
      scriptUtxo.input.txHash,
      scriptUtxo.input.outputIndex,
      scriptUtxo.output.amount,
      scriptUtxo.output.address
    )
    .txInScript(scriptCbor)
    .txInRedeemerValue(redeemer)
    .txInInlineDatumPresent()
    .txInCollateral(
      collateral.input.txHash,
      collateral.input.outputIndex,
      collateral.output.amount,
      collateral.output.address
    )
    .changeAddress(walletAddress)
    .selectUtxosFrom(utxos)
    .complete();

  const signedTx = await wallet.signTx(txBuilder.txHex);
  const txHash = await wallet.submitTx(signedTx);
  console.log("unlocked at:", txHash);
}

main();
\end{verbatim}

Three details deserve attention.

\begin{itemize}
  \item \texttt{txInInlineDatumPresent()} tells the builder the datum is already
        on-chain. If the lock used a datum hash you would call
        \texttt{txInDatumValue(datum)} and supply the preimage instead.
  \item \texttt{txInCollateral} is mandatory for any Plutus spend. If the script
        fails, the collateral is forfeited, which is why you test with
        \texttt{aiken check} first, then on preprod, then on mainnet.
  \item There is no \texttt{requiredSignerHash}. Our validator does not check
        signatures at all: possession of a valid proof \emph{is} the
        authorisation. That is the whole point, and also, as
        Chapter~\ref{ch:password} warned, why replay protection matters
        (Section~\ref{sec:ch7-ops}).
\end{itemize}

\subsection*{Running It}

\begin{verbatim}
export BLOCKFROST_PROJECT_ID=preprodXXXXXXXXXXXXXXXX
npx tsx offchain/lock.ts
npx tsx offchain/unlock.ts <lock-tx-hash>
\end{verbatim}

%--------------------------------------------------------------------------
\section{What the Transactions Actually Cost} \label{sec:ch7-cost}

Abstract talk about efficiency is worth less than measurement. Aiken reports
execution units for every test, so we can read the real cost of our verifier
directly:

\begin{verbatim}
aiken check
\end{verbatim}

\begin{verbatim}
  PASS  zk_password/password_tests :: groth16_verify_valid_proof
        mem: 34219, cpu: 2005373690
\end{verbatim}

Compare against the live mainnet protocol parameters (epoch 646, Conway era,
protocol version 11):

\begin{center}
\begin{tabular}{lrr}
  \hline
  \textbf{Resource} & \textbf{Used} & \textbf{Budget} \\
  \hline
  CPU steps            & 2{,}005{,}373{,}690 & 10{,}000{,}000{,}000 \\
  Memory units         & 34{,}219            & 16{,}500{,}000 \\
  Script size (bytes)  & 570                 & 16{,}384 (whole tx) \\
  \hline
\end{tabular}
\end{center}

So a Groth16 verification consumes \textbf{20.1\% of the CPU budget} and
\textbf{0.21\% of the memory budget}. At the current prices, 0.0000721
lovelace per CPU step and 0.0577 lovelace per memory unit, the script
execution costs:

\[
  2{,}005{,}373{,}690 \times 0.0000721 \;+\; 34{,}219 \times 0.0577
  \;\approx\; 146{,}562 \text{ lovelace} \;\approx\; \mathbf{0.147\ ADA}
\]

\begin{highlight}
\textbf{Verifying a zero-knowledge proof on Cardano costs about 0.15 ADA and
one fifth of a transaction's CPU budget, and that figure does not depend on
how complicated the circuit was.} Groth16 verification is three point
decompressions and four pairings, whether the statement had 216 constraints or
216{,}000.
\end{highlight}

This constancy is the single most important economic fact about ZK on Cardano,
and we will lean on it heavily in Chapter~\ref{ch:tornado}. Its immediate
corollary is a hard ceiling:

\[
  \left\lfloor \frac{10{,}000{,}000{,}000}{2{,}005{,}373{,}690} \right\rfloor
  = 4 \text{ Groth16 verifications per transaction}
\]

You cannot batch five ZK spends into one transaction. The memory budget would
allow 482 of them; the CPU budget allows four.

%--------------------------------------------------------------------------
\section{Generating the Off-Chain Code: \texttt{aiken-zk}} \label{sec:ch7-aikenzk}

Writing all of this by hand, once, is educational. Writing it for every circuit
is toil. The \textbf{aiken-zk} tool from
Eryx\footnote{\texttt{https://github.com/eryxcoop/cardano-zk-aiken}}, funded
by Project Catalyst Fund 13, automates most of it.

The idea is a language extension: you annotate an Aiken validator with an
\texttt{offchain} block describing a computation, and the tool emits the Circom
circuit, the standard Aiken validator containing the Groth16 verifier, and the
off-chain proving glue.

\subsection*{The Workflow}

\begin{verbatim}
aiken-zk new my_project

# transform annotated Aiken into plain Aiken + circuit + vk
aiken-zk build code_with_offchain.ak aiken_code.ak

# produce a proof as an Aiken literal (for tests)
aiken-zk prove aiken output.circom verification_key.zkey \
    inputs.json proof.ak

# produce a proof as a ready-made MeshJS redeemer
aiken-zk prove meshjs output.circom verification_key.zkey \
    inputs.json zk_redeemer.ts
\end{verbatim}

That last command is the interesting one for this chapter: it emits the
TypeScript redeemer object directly, replacing Stage 6 and the hand-written
\texttt{mConStr0} nesting of Section~\ref{sec:ch7-encoding}.

\subsection*{The \texttt{offchain} Vocabulary}

\begin{center}
\begin{tabular}{ll}
  \hline
  \textbf{Statement} & \textbf{Proves} \\
  \hline
  \texttt{addition(x, y, z)}        & $x + y = z$ \\
  \texttt{subtraction(x, y, z)}     & $x - y = z$ \\
  \texttt{multiplication(x, y, z)}  & $x \cdot y = z$ \\
  \texttt{fibonacci(x, y, t, z)}    & $t$-th Fibonacci term is $z$ \\
  \texttt{if(x, y, z, w)}           & $y = z$ if $x=1$, else $y = w$ \\
  \texttt{assert\_eq(x, y)}         & $x = y$ \\
  \texttt{sha256(t, A, B)}          & $\mathrm{sha256}(A) = B$ \\
  \texttt{poseidon(t, A, x)}        & $\mathrm{Poseidon}(A) = x$ \\
  \texttt{polynomial\_evaluations}  & polynomial evaluation \\
  \texttt{merkle\_tree\_checker}    & Merkle inclusion path \\
  \texttt{custom("...circom", [..])}& an arbitrary Circom component \\
  \hline
\end{tabular}
\end{center}

Arguments carry visibility modifiers, so \texttt{offchain addition(priv x, pub y, z)}
makes \texttt{x} a private witness and \texttt{y} a public input. The
\texttt{custom} escape hatch is what makes the tool genuinely useful: you can
drop in any Circom component and still get the verifier and redeemer generated.

\subsection*{Current Limitations}

The project is young, and its README is candid about the edges:

\begin{itemize}
  \item Only \textbf{one} \texttt{offchain} statement per Aiken file.
  \item \texttt{fibonacci}'s length parameter must be a literal.
  \item List arguments cannot be literals.
  \item \textbf{The \texttt{merkle\_tree\_checker} hash function is mocked:
        ``do not use in production''.}
\end{itemize}

That last bullet is not an incidental bug. It is the same wall we will hit
throughout Chapter~\ref{ch:tornado}, and it has a specific cause:
Cardano has no ZK-friendly hash primitive. We will measure exactly how
expensive that absence is.

\begin{remark}
Use \texttt{aiken-zk} to bootstrap and to learn the shape of the generated
code, but read what it emits. A tool that hides the Circom layer will
eventually hide a curve mismatch or a public-input ordering bug from you, and
those are precisely the bugs that cost collateral.
\end{remark}

%--------------------------------------------------------------------------
\section{Operational Concerns} \label{sec:ch7-ops}

A working pipeline is not a working product. Five issues bite in production.

\subsection*{1. Replay Protection}

Our validator accepts any valid proof for the hash in the datum. A proof is
public data sitting in a submitted transaction. An observer watching the mempool
can copy the proof and submit their own transaction paying themselves.

The fix is to bind the proof to something unique. Add a public input the
validator can independently check:

\begin{itemize}
  \item \textbf{A nonce in the datum.} The validator requires the continuing
        output's datum to carry \texttt{nonce + 1}, and the circuit takes the
        nonce as a public input. Janus Wallet uses exactly this, with
        \texttt{challenge} and \texttt{challengeFlag} signals.
  \item \textbf{The recipient address.} Make the payment destination a public
        input. A stolen proof can then only pay the original recipient, which
        removes the attacker's incentive.
  \item \textbf{The spent output reference.} Bind the proof to the specific UTxO
        being consumed, which is unique by construction.
\end{itemize}

\begin{highlight}
Any ZK validator that authorises value transfer on the strength of a proof
alone, with no binding to the transaction, is replayable. This is not a Cardano
quirk; it is true on every chain. Treat a proof as a bearer instrument.
\end{highlight}

\subsection*{2. Where to Put the Verification Key}

Three options, with real trade-offs:

\begin{center}
\begin{tabular}{lll}
  \hline
  \textbf{Placement} & \textbf{Cost} & \textbf{Upgradeable} \\
  \hline
  Inline datum          & $\sim$430 bytes per UTxO & Yes, per UTxO \\
  Compile-time param    & Larger script, once      & No (new address) \\
  Reference input       & 15 lovelace/byte fee     & Yes, centrally \\
  \hline
\end{tabular}
\end{center}

Baking the VK into the script with \texttt{applyParamsToScript} is the cheapest
per transaction and the most rigid: changing the circuit changes the script
hash, hence the address, and funds at the old address need migrating. This is
the trap Janus Wallet documents, a password change migrates the address, and
funds sent to the old one are unrecoverable.

The reference-input approach stores the VK in a UTxO consulted read-only by many
transactions. With \texttt{minFeeRefScriptCoinsPerByte} currently at 15, our
570-byte validator costs $570 \times 15 = 8{,}550$ lovelace ($\approx$ 0.009
ADA) per transaction when referenced, cheaper than including the script
inline once you have more than a handful of spends.

\subsection*{3. Distributing the Ceremony Artifacts}

Every prover needs \texttt{circuit.wasm} and \texttt{zkey\_final.zkey}. These
are not secrets, but they are integrity-critical: a substituted zkey lets an
attacker forge proofs that your on-chain VK will happily accept only if the VK
was substituted too, so pin both. Publish the hashes, ideally on-chain
alongside the VK, and have clients check them.

\subsection*{4. Budget Estimation and Collateral}

Let the provider's evaluator compute execution units rather than hard-coding
them. Hard-coded budgets from a tutorial will be wrong after the next cost-model
update, and an underestimate means a failed script and forfeited collateral.
Note that \texttt{maxCollateralInputs} is 3 and
\texttt{collateralPercentage} is 150, so you need a clean, ADA-only UTxO of
sufficient size available at all times.

\subsection*{5. Indexing}

Nothing in Cardano tells you ``here are the UTxOs at my script address with a
datum matching this hash''. You need to query and filter. Blockfrost's
\texttt{/addresses/\{addr\}/utxos} plus datum decoding works for small
applications; anything larger wants Kupo, Ogmios, or Carp maintaining an index.
Chapter~\ref{ch:tornado} shows what happens when the off-chain indexer must
reconstruct cryptographic state, not just find UTxOs.

%--------------------------------------------------------------------------
\section{Failure Modes and How to Recognise Them} \label{sec:ch7-failures}

Almost every off-chain ZK bug on Cardano presents as the same unhelpful
symptom. This table is the debugging shortcut this chapter exists to provide.

\begin{center}
\begin{tabular}{p{4.6cm}p{7.4cm}}
  \hline
  \textbf{Symptom} & \textbf{Cause} \\
  \hline
  snarkjs verifies, chain rejects
    & Circuit compiled for BN254. Recompile with
      \texttt{--prime bls12381}. \\
  \addlinespace
  \texttt{decompress} fails / script errors immediately
    & Uncompressed points in datum or redeemer. Run
      \texttt{compress\_vk.js} and \texttt{compress\_proof.js}. \\
  \addlinespace
  Constraint unsatisfiable during \texttt{fullprove}
    & Hash computed with \texttt{circomlibjs} (BN254) instead of the
      BLS12-381 helper circuit. \\
  \addlinespace
  \texttt{I caught an opaque type trying to escape}
    & \texttt{Scalar} in a datum or redeemer. Use \texttt{Int} and
      \texttt{scalar.new} inside the validator. \\
  \addlinespace
  Verification fails with a correct-looking hash
    & \texttt{parseInt}/\texttt{Number} truncated the 255-bit field
      element. Use \texttt{BigInt}. \\
  \addlinespace
  Fails only for some public inputs
    & \texttt{vk\_ic} order wrong. $IC_0$ is the constant term;
      $IC_{i+1}$ pairs with public input $i$. \\
  \addlinespace
  \texttt{ExUnits} exceeded
    & More than four Groth16 verifications, or on-chain field
      arithmetic. See Chapter~\ref{ch:tornado}. \\
  \addlinespace
  Datum mismatch at spend time
    & Field order differs between the Aiken \texttt{type} and the
      \texttt{mConStr0} array. \\
  \hline
\end{tabular}
\end{center}

\begin{remark}
Notice how many of these are \emph{encoding} bugs rather than cryptographic
ones. The mathematics is handled by libraries that are correct. The bugs live in
the seams: field choice, point format, integer width, field order. Build a
round-trip test for your encoding before you build a transaction.
\end{remark}

%--------------------------------------------------------------------------
\section{Chapter Summary} \label{sec:ch7-summary}

\begin{itemize}
  \item The on-chain component of a ZK application on Cardano is tiny: 570
        bytes of Plutus Core for a complete Groth16 verifier.
  \item Circom + snarkjs + Groth16 over BLS12-381 is the only fully supported
        path today. \texttt{--prime bls12381} is not optional.
  \item Six pipeline stages: three per circuit (compile, setup, compress VK),
        three per proof (witness, prove, compress proof).
  \item Where the prover runs decides whether the system is actually private.
        Client-side proving is the only architecture that delivers
        zero-knowledge against the operator.
  \item Plutus Data encoding is positional. \texttt{BigInt}, correct field
        order, compressed points.
  \item MeshJS builds both transactions in about thirty lines each:
        \texttt{txOutInlineDatumValue} to lock,
        \texttt{spendingPlutusScript("V3")} with
        \texttt{txInRedeemerValue} to unlock.
  \item A measured Groth16 verification costs 2{,}005{,}373{,}690 CPU steps and
        34{,}219 memory units, 20\% of the CPU budget, $\approx$ 0.147 ADA,
        and \emph{constant} in circuit size. Four verifications per transaction
        is the hard ceiling.
  \item \texttt{aiken-zk} generates most of this from an \texttt{offchain}
        annotation, including the MeshJS redeemer, with the notable caveat
        that its Merkle tree hash is mocked.
  \item Proofs are bearer instruments. Bind them to a nonce, a recipient, or an
        output reference, or they are replayable.
\end{itemize}

The next chapter takes everything in this one and points it at a hard target: a
Tornado Cash--style mixer. We will find that the verification side works fine,
and that the problem lies somewhere else entirely.
